% !TEX program = xelatex
\documentclass[11pt,a4paper]{article}
\usepackage{algo-preamble}

\title{\textbf{L09 \textperiodcentered{} Γραφήματα IV}\\[2pt]
\large Dijkstra \& Ελάχιστα Συνδετικά Δέντρα}
\author{Algorithms Class Hub}
\date{\small Αλγόριθμοι και Πολυπλοκότητα (K17) \textperiodcentered{} ΕΚΠΑ}

\begin{document}
\maketitle

\begin{center}\itshape\small
Συντομότερες διαδρομές με βάρη (Dijkstra), ελάχιστα συνδετικά δέντρα (MST), η
ιδιότητα αποκοπής και η ιδιότητα κύκλου, και οι άπληστοι αλγόριθμοι Prim και
Kruskal.
\end{center}

\section{Τι θα δούμε}

Δύο μεγάλα προβλήματα σε ζυγισμένα γραφήματα, και τρεις \textbf{άπληστους}
αλγορίθμους που τα λύνουν:
\begin{enumerate}
  \item \textbf{Συντομότερη διαδρομή με θετικά βάρη} $\to$ \textbf{Dijkstra},
  $O(m \log n)$.
  \item \textbf{Ελάχιστο Συνδετικό Δέντρο (ΕΕΔ / MST)} $\to$ \textbf{Prim},
  \textbf{Kruskal} (και η «αντίστροφη διαγραφή»), $O(m \log n)$.
\end{enumerate}

Και οι τρεις είναι \textbf{άπληστοι}: σε κάθε βήμα παίρνουν την τοπικά καλύτερη
απόφαση. Το ότι αυτό δίνει \textbf{παγκόσμια} βέλτιστη λύση δεν είναι προφανές
--- θα δούμε την απόδειξη μέσω δύο όμορφων ιδιοτήτων: της \textbf{ιδιότητας
αποκοπής} και της \textbf{ιδιότητας κύκλου}.

\begin{keypoint}
\textbf{Σπουδαία διάλεξη για την εξεταστική.} Σχεδόν κάθε χρονιά υπάρχει θέμα που
ζητάει «τρέξε Dijkstra στο γράφημα\dots», «βρες το ΕΕΔ\dots», ή «εφάρμοσε την
ιδιότητα αποκοπής για να αποδείξεις\dots». Επίσης είναι τα τελευταία θέματα όπου
το BFS/DFS από προηγούμενα μαθήματα παίζει ως υπορουτίνα.
\end{keypoint}

\section{Συντομότερη Διαδρομή με Βάρη}

\subsection{Διατύπωση}

\textbf{Είσοδος:} κατευθυνόμενο γράφημα $G = (V, A)$. Σε κάθε ακμή $e \in A$
ανατίθεται \textbf{θετικός} πραγματικός αριθμός $\ell_e > 0$ (το «μήκος» της).

\textbf{Κόστος διαδρομής:} το άθροισμα των μηκών των ακμών της.

\textbf{Πρόβλημα:} δοθέντων αφετηρίας $s$ και προορισμού $t$, βρες την
κατευθυνόμενη διαδρομή από $s$ σε $t$ με το \textbf{ελάχιστο} συνολικό μήκος.

\subsection{Γιατί δεν αρκεί το BFS}

Από το προηγούμενο μάθημα ξέρουμε ότι το BFS λύνει το πρόβλημα \textbf{χωρίς}
βάρη. Σε ζυγισμένο γράφημα όμως, η διαδρομή με τις λιγότερες ακμές \textbf{δεν}
είναι κατ' ανάγκη η μικρότερη σε κόστος: μια διαδρομή 3 ακμών μήκους 5 η καθεμία
(σύνολο 15) είναι \textbf{καλύτερη} από μια διαδρομή 1 ακμής μήκους 100.

\section{Αλγόριθμος Dijkstra}

\subsection{Ιδέα --- άπληστη επέκταση από την $s$}

Κρατάμε ένα σύνολο $S$ από κορυφές \textbf{εξερευνημένες}, για τις οποίες έχουμε
ήδη υπολογίσει την απόσταση από την $s$. Σε κάθε βήμα:
\begin{enumerate}
  \item Για κάθε \textbf{μη εξερευνημένη} κορυφή $v$, υπολογίζουμε
  \[
  \pi(v) = \min_{u \in S}\bigl\{ d(u) + \ell_{u,v} \bigr\}
  \]
  δηλαδή «αν φτάσεις στην $v$ από κάποια $u \in S$, ποιο είναι το ελάχιστο
  κόστος;»
  \item Διάλεξε την κορυφή $v$ με το \textbf{ελάχιστο} $\pi(v)$.
  \item Θέσε $d(v) := \pi(v)$, πρόσθεσε την $v$ στο $S$.
\end{enumerate}
\textbf{Αρχικοποίηση:} $S = \{s\}$, $d(s) = 0$.

\begin{intuition}
\textbf{Διαισθητικά:} φαντάσου ότι τοποθετείς ένα μαθηματικό «κύμα» στην $s$ που
διαδίδεται με ταχύτητα 1. Σε κάθε χρονική στιγμή, η επόμενη κορυφή που το κύμα
φτάνει είναι ακριβώς η μη-εξερευνημένη με το μικρότερο $\pi(v)$. Είναι μια
\textbf{γενίκευση του BFS} όπου το «επόμενο επίπεδο» δεν είναι μια ακμή μακριά,
αλλά \textbf{όσο επιτρέπει η μικρότερη ακμή}.
\end{intuition}

\subsection{Ψευδοκώδικας}

\begin{codebox}
\begin{verbatim}
Dijkstra(G, l, s):
  Για κάθε v ∈ V:
    d[v] <- ∞
    π[v] <- ∅
  d[s] <- 0
  PQ <- ουρά προτεραιότητας με όλες τις κορυφές (κλειδί = d)
  S <- ∅
  Ενόσω PQ ≠ ∅:
    u <- extractMin(PQ)
    S <- S ∪ {u}
    Για κάθε ακμή (u, v) ∈ A:
      Εάν v ∉ S και d[u] + l(u,v) < d[v]:
        d[v] <- d[u] + l(u,v)
        π[v] <- u
        decreaseKey(PQ, v, d[v])
  επίστρεψε d, π
\end{verbatim}
\end{codebox}

Όταν τελειώσει, $d[v]$ είναι η συντομότερη απόσταση από την $s$ στην $v$. Με
τους «γονείς» $\pi[\cdot]$ μπορούμε να ανακατασκευάσουμε την ίδια τη διαδρομή ---
ακολουθώντας πίσω από $t$ προς $s$.

\subsection{Ορθότητα --- αναλλοίωτη με επαγωγή στο $|S|$}

\begin{keypoint}
\textbf{Αναλλοίωτη:} για κάθε $u \in S$, $d(u)$ ισούται με το μήκος της
\textbf{πραγματικής} συντομότερης διαδρομής $s \to u$.
\end{keypoint}

\textbf{Βάση:} $|S| = 1$, $S = \{s\}$, $d(s) = 0$. $\checkmark$

\textbf{Επαγωγικό βήμα:} υποθέτουμε ότι η αναλλοίωτη ισχύει για $|S| = k$. Έστω
$v$ η κορυφή που μπαίνει στο $S$ στο βήμα $k+1$ --- αυτή που ελαχιστοποιεί το
$\pi(\cdot)$ --- και έστω $(u, v)$ η ακμή που δίνει αυτό το ελάχιστο, δηλαδή
$\pi(v) = d(u) + \ell_{u,v}$.

Θα δείξουμε ότι \textbf{κάθε} άλλη διαδρομή $P : s \to v$ έχει μήκος
$\ge \pi(v)$. Έστω λοιπόν $P$ μια τέτοια διαδρομή. Έστω $(x, y)$ η
\textbf{πρώτη} ακμή της $P$ που «βγαίνει» από το $S$ (δηλαδή $x \in S$,
$y \notin S$). Έστω $P'$ η υποδιαδρομή $s \to x$:
\[
\begin{aligned}
\ell(P) &\ge \ell(P') + \ell_{x,y} && \text{(αρχικό κομμάτι της } P\text{)} \\
        &\ge d(x) + \ell_{x,y}     && \text{(αναλλοίωτη επαγωγής για } x \in S
\text{)} \\
        &\ge \pi(y)                && \text{(από τον ορισμό του } \pi(y)
\text{)} \\
        &\ge \pi(v)                && \text{(ο Dijkstra διάλεξε } v \text{ ως
minimum)}
\end{aligned}
\]
Άρα $\ell(P) \ge \pi(v) = d(v)$. $\blacksquare$

\begin{warningbox}
\textbf{Σημαντικός περιορισμός:} η αναλλοίωτη χρησιμοποιεί ότι \textbf{όλα τα
βάρη είναι θετικά}. Αν υπήρχε αρνητική ακμή, θα μπορούσε να «μειώσει» το κόστος
μιας υποδιαδρομής που βγαίνει και ξαναμπαίνει στο $S$ --- άκυρη απόδειξη, και
πράγματι ο Dijkstra \textbf{δεν} δουλεύει με αρνητικά βάρη. Σε εκείνη την
περίπτωση χρειάζεται \textbf{Bellman-Ford} ($O(n m)$).
\end{warningbox}

\subsection{Πολυπλοκότητα}

\begin{itemize}
  \item \textbf{Χωρίς δομή δεδομένων:} σε κάθε βήμα ψάχνουμε γραμμικά την κορυφή
  με το μικρότερο $\pi$. $O(n)$ ανά extraction, $n$ extractions $\to O(n^2)$.
  Καλό για πυκνά γραφήματα.
  \item \textbf{Με δυαδικό σωρό (heap):} ένας σωρός υποστηρίζει \texttt{extractMin}
  σε $O(\log n)$ και \texttt{decreaseKey} σε $O(\log n)$. Συνολικά: $n$
  extractMin $\to O(n \log n)$· μέχρι $m$ decreaseKey $\to O(m \log n)$.
  \item \textbf{Σύνολο:} $O((n + m) \log n) = O(m \log n)$ (για συνεκτικά
  γραφήματα όπου $m \ge n - 1$).
\end{itemize}
Με Fibonacci heap πέφτει σε $O(m + n \log n)$ --- αλλά εκείνο πέρα από το μάθημα.

\section{Ελάχιστο Συνδετικό Δέντρο (ΕΕΔ / MST)}

\subsection{Διατύπωση}

\textbf{Είσοδος:} μη κατευθυνόμενο \textbf{συνεκτικό} γράφημα $G = (V, E)$ με
κόστη $c_e$ σε κάθε ακμή.

\textbf{Επικαλύπτον δέντρο} (spanning tree): υποσύνολο $T \subseteq E$ τέτοιο
ώστε (1) το $T$ είναι \textbf{δέντρο} (συνεκτικό $+$ ακυκλικό) και (2)
επικαλύπτει \textbf{όλες} τις κορυφές $V$. Από προηγούμενο μάθημα ξέρουμε ότι
κάθε δέντρο με $n$ κορυφές έχει $n - 1$ ακμές.

\textbf{Πρόβλημα:} βρες ένα επικαλύπτον δέντρο \textbf{ελάχιστου κόστους},
δηλαδή ελαχιστοποίησε $\sum_{e \in T} c_e$.

\subsection{Γιατί δεν αρκεί η ωμή βία}

\begin{keypoint}
\textbf{Θεώρημα του Cayley.} Το πλήρες γράφημα $K_n$ έχει \textbf{$n^{n-2}$}
διαφορετικά επικαλύπτοντα δέντρα.
\end{keypoint}

Άρα η εξαντλητική απαρίθμηση είναι \textbf{ασύμφορη} --- π.χ.\ το $K_{10}$ έχει
100 εκατομμύρια. Χρειαζόμαστε έξυπνο αλγόριθμο.

\textbf{Πραγματικές εφαρμογές:} σχεδίαση δικτύων (τηλεφωνικά, ηλεκτρικά,
υδρευτικά, καλωδιακά, οδικά), προσεγγιστικοί αλγόριθμοι για NP-δύσκολα προβλήματα
(TSP, δέντρο Steiner), max-bottleneck paths, LDPC κώδικες για επιδιόρθωση
λαθών, καταχώρηση εικόνας με εντροπία Rényi, το spanning tree protocol σε δίκτυα
Ethernet (αποφυγή loops), και ανάλυση συστάδων (clustering).

\section{Τρεις Άπληστοι Αλγόριθμοι}

\begin{center}
\begin{tabular}{@{}l p{6.6cm} l@{}}
\toprule
Αλγόριθμος & Στρατηγική & Πολυπλοκότητα \\
\midrule
\textbf{Kruskal} & ακμές σε \textbf{αύξουσα} σειρά κόστους, προσθέτεις αν δεν
δημιουργεί κύκλο & $O(m \log n)$ \\
\textbf{Αντίστροφη διαγραφή} & ακμές σε \textbf{φθίνουσα} σειρά κόστους, αφαιρείς
αν δεν αποσυνδέει το γράφημα & $O(m \log n)$ \\
\textbf{Prim} & ξεκίνα από κορυφή $s$, σε κάθε βήμα πρόσθεσε τη φθηνότερη ακμή
που «μεγαλώνει» το δέντρο & $O(m \log n)$ \\
\bottomrule
\end{tabular}
\end{center}

\textbf{Όλοι} δίνουν την ίδια λύση. Διαφορετική διατύπωση, ίδιο αποτέλεσμα.

\section{Κύκλοι, Αποκοπές, και η Τομή τους}

Πριν αποδείξουμε ότι οι παραπάνω αλγόριθμοι είναι σωστοί, χρειαζόμαστε δύο
ορισμούς.

\subsection{Αποκοπή}

\textbf{Αποκοπή} (cut): οποιοδήποτε υποσύνολο κορυφών $A \subset V$. Το
αντίστοιχο \textbf{σύνολο ακμών αποκοπής} $D$ είναι το σύνολο των ακμών που έχουν
\textbf{το ένα άκρο στο $A$} και το άλλο στο $V \setminus A$.

\begin{intuition}
\textbf{Διαισθητικά:} η αποκοπή «κόβει» το γράφημα σε δύο πλευρές. Οι ακμές
αποκοπής είναι αυτές που «περνάνε» τη γραμμή.
\end{intuition}

\subsection{Το βασικό λήμμα}

\begin{keypoint}
\textbf{Λήμμα (κύκλος-αποκοπή).} Ένας κύκλος $C$ και ένα σύνολο ακμών αποκοπής
$D$ τέμνονται σε \textbf{άρτιο} πλήθος ακμών.

\textbf{Απόδειξη:} αν ο κύκλος είναι ολόκληρος μέσα στο $A$ ή ολόκληρος έξω, η
τομή είναι κενή (= 0). Αλλιώς, κάθε φορά που ο κύκλος «μπαίνει» στο $A$ (μέσω
ακμής του $D$), πρέπει επίσης να «βγει» (επίσης μέσω ακμής του $D$). Άρα οι
ακμές του $C$ στο $D$ ζευγαρώνουν $\to$ άρτιο πλήθος. $\blacksquare$
\end{keypoint}

\subsection{Ιδιότητα Αποκοπής}

\begin{keypoint}
\textbf{Ιδιότητα Αποκοπής.} Έστω $S \subset V$ οποιαδήποτε αποκοπή και έστω $e$
η ακμή \textbf{ελάχιστου κόστους} στο σύνολο ακμών αποκοπής $D(S)$. Τότε
\textbf{το ΕΕΔ περιέχει την $e$}.
\end{keypoint}

\textbf{Απόδειξη (με ανταλλαγή).} Έστω $T^*$ το ΕΕΔ, και υποθέτουμε για αντίφαση
ότι $e \notin T^*$.
\begin{enumerate}
  \item Προσθέτοντας την $e$ στο $T^*$ δημιουργείται μοναδικός κύκλος $C$
  (δέντρο $+$ ακμή $=$ ένας κύκλος).
  \item Η $e$ ανήκει στον $C$ \textbf{και} στο $D(S)$.
  \item Από το λήμμα κύκλου-αποκοπής, υπάρχει \textbf{τουλάχιστον μία άλλη} ακμή
  $f \in C \cap D(S)$, $f \ne e$.
  \item Φτιάχνουμε νέο δέντρο: $T' = T^* \cup \{e\} \setminus \{f\}$. Είναι
  επικαλύπτον δέντρο (αφαίρεση μιας ακμής κύκλου διατηρεί τη συνεκτικότητα).
  \item Αφού $c_e < c_f$ (η $e$ είναι η ελάχιστη στο $D$), έχουμε
  $\text{cost}(T') < \text{cost}(T^*)$ --- \textbf{αντίφαση} ότι το $T^*$ ήταν
  ελάχιστο. $\blacksquare$
\end{enumerate}

\textbf{Εφαρμογή:} Prim και Kruskal εξασφαλίζουν ότι κάθε ακμή που προσθέτουν
είναι ελάχιστη σε κάποια αποκοπή $\to$ από την ιδιότητα ανήκει στο ΕΕΔ.

\subsection{Ιδιότητα Κύκλου}

\begin{keypoint}
\textbf{Ιδιότητα Κύκλου.} Έστω $C$ οποιοσδήποτε κύκλος και έστω $f$ η ακμή
\textbf{μέγιστου κόστους} στο $C$. Τότε \textbf{το ΕΕΔ δεν περιέχει την $f$}.
\end{keypoint}

\textbf{Απόδειξη (συμμετρική).} Υποθέτουμε για αντίφαση ότι $f \in T^*$.
\begin{enumerate}
  \item Αφαιρώντας την $f$ από το $T^*$ δημιουργείται αποκοπή $S$ (δέντρο
  $-$ ακμή $=$ δύο συνιστώσες).
  \item Η $f$ ανήκει στον $C$ \textbf{και} στο $D(S)$.
  \item Από το λήμμα, υπάρχει $e \in C \cap D(S)$, $e \ne f$.
  \item $T' = T^* \cup \{e\} \setminus \{f\}$ είναι επικαλύπτον δέντρο.
  \item Αφού $c_e < c_f$ (η $f$ είναι η μέγιστη στο $C$),
  $\text{cost}(T') < \text{cost}(T^*)$. \textbf{Αντίφαση.} $\blacksquare$
\end{enumerate}

\textbf{Εφαρμογή:} αντίστροφη διαγραφή και Kruskal εξασφαλίζουν ότι κάθε ακμή που
\textbf{αποκλείουν} είναι μέγιστη σε κάποιο κύκλο.

\begin{intuition}
\textbf{Παραδοχή απλοποίησης:} στις αποδείξεις υποθέσαμε ότι \textbf{όλα τα κόστη
είναι μοναδικά} (κανένα ζεύγος ακμών δεν έχει ίδιο κόστος). Όταν υπάρχουν
ισοτιμίες, οι αλγόριθμοι παίρνουν αυθαίρετες αποφάσεις και υπάρχουν πολλά
«ελάχιστα» δέντρα. Όλα έχουν το ίδιο \textbf{συνολικό} κόστος --- αλλά
διαφορετικές συγκεκριμένες ακμές.
\end{intuition}

\section{Αλγόριθμος του Prim}

\textbf{Ιδέα.} Ξεκινάμε με μία κορυφή και «μεγαλώνουμε» το δέντρο \textbf{προς τα
έξω}. Σε κάθε βήμα, η αποκοπή είναι «κορυφές μέσα στο δέντρο» vs «κορυφές έξω».
Από την \textbf{ιδιότητα αποκοπής}, η φθηνότερη ακμή που περνάει αυτή τη γραμμή
ανήκει στο ΕΕΔ --- οπότε την παίρνουμε.

\begin{codebox}
\begin{verbatim}
Prim(G, c, s):
  Για κάθε v ∈ V:
    α[v] <- ∞               // κόστος επισύναψης της v στο S
    π[v] <- ∅
  α[s] <- 0
  Q <- ουρά προτεραιότητας με όλες τις κορυφές (κλειδί = α)
  S <- ∅
  T <- ∅
  Ενόσω Q ≠ ∅:
    u <- extractMin(Q)
    S <- S ∪ {u}
    Εάν π[u] ≠ ∅:
      T <- T ∪ { (u, π[u]) }
    Για κάθε ακμή e = {u, v} του G:
      Εάν v ∉ S και c_e < α[v]:
        α[v] <- c_e
        π[v] <- u
        decreaseKey(Q, v, α[v])
  επίστρεψε T
\end{verbatim}
\end{codebox}

\textbf{Παρατήρηση:} μοιάζει εξαιρετικά με Dijkstra. Η διαφορά είναι ότι το
\textbf{κλειδί} $\alpha[v]$ δεν είναι «κόστος διαδρομής από $s$» αλλά «κόστος της
ελάχιστης ακμής που συνδέει τη $v$ με το δέντρο». Άρα ο Prim είναι
\textbf{τοπικός} (κοιτάει μόνο μία ακμή), ενώ ο Dijkstra είναι
\textbf{συσσωρευτικός} (αθροίζει την πορεία).

\textbf{Πολυπλοκότητα.} Ίδια ανάλυση με Dijkstra: $O(m \log n)$ με δυαδικό σωρό,
$O(n^2)$ με απλό πίνακα.

\section{Αλγόριθμος του Kruskal}

\textbf{Ιδέα.} Ταξινομούμε όλες τις ακμές σε \textbf{αύξουσα} σειρά κόστους.
Σαρώνουμε με τη σειρά:
\begin{itemize}
  \item \textbf{Περίπτωση 1:} αν η ακμή $e = \{u, v\}$ δημιουργεί κύκλο στο
  τρέχον $T$ $\to$ απορρίπτουμε (ιδιότητα κύκλου: η $e$ είναι η μεγαλύτερη ή ίση
  στον κύκλο, αφού όλες οι ήδη μπασμένες είναι φθηνότερες).
  \item \textbf{Περίπτωση 2:} αλλιώς $\to$ προσθέτουμε στο $T$ (ιδιότητα
  αποκοπής: αν δεν δημιουργεί κύκλο, ενώνει δύο συνεκτικές συνιστώσες, και η $e$
  είναι η ελάχιστη στην αποκοπή «μία συνιστώσα vs την άλλη»).
\end{itemize}
Σταματάμε όταν το $T$ έχει $n - 1$ ακμές.

\begin{codebox}
\begin{verbatim}
Kruskal(G, c):
  Ταξινόμησε τις ακμές: c_1 <= c_2 <= ... <= c_m
  T <- ∅
  Για κάθε u ∈ V:
    MakeSet(u)               // κάθε κορυφή ξεκινά σε δικό της σύνολο
  Για i = 1 έως m:
    {u, v} = e_i
    Εάν Find(u) ≠ Find(v):   // όχι στην ίδια συνιστώσα -> όχι κύκλος
      T <- T ∪ {e_i}
      Union(u, v)
  επίστρεψε T
\end{verbatim}
\end{codebox}

Ο έλεγχος «οι $u, v$ είναι ήδη συνδεδεμένες;» γίνεται με τη δομή
\textbf{Union-Find} που θα δούμε αναλυτικά σε επόμενο μάθημα (δομές δεδομένων).

\textbf{Πολυπλοκότητα:}
\begin{itemize}
  \item Ταξινόμηση ακμών: $O(m \log m) = O(m \log n)$.
  \item $m$ κλήσεις σε \texttt{Find}/\texttt{Union}: $O(m \cdot \alpha(n))$ με
  union-by-rank $+$ path compression, όπου $\alpha$ είναι η \textbf{αντίστροφη
  Ackermann} --- πρακτικά σταθερά.
  \item \textbf{Σύνολο:} $O(m \log n)$ (η ταξινόμηση είναι το κυρίαρχο κόστος).
\end{itemize}

\section{Σύγκριση Prim vs Kruskal}

\begin{center}
\begin{tabular}{@{}l p{4.6cm} p{4.6cm}@{}}
\toprule
 & \textbf{Prim} & \textbf{Kruskal} \\
\midrule
Δομή που μεγαλώνει & ένα δέντρο που μεγαλώνει από κορυφή & δάσος από πολλά
δέντρα που συγχωνεύονται \\
Δομή δεδομένων & priority queue πάνω σε κορυφές & sort $+$ Union-Find \\
Πολυπλοκότητα & $O(m \log n)$ & $O(m \log n)$ \\
Καταλληλότερος για & πυκνά γραφήματα & αραιά γραφήματα, παράλληλη υλοποίηση \\
\bottomrule
\end{tabular}
\end{center}

\textbf{Στην εξεταστική:} συχνά ζητείται να εφαρμοστούν και οι δύο σε ένα
συγκεκριμένο γράφημα. Δίνουν την ίδια λύση (αν τα κόστη είναι μοναδικά), αλλά
\textbf{η σειρά} που προσθέτουν τις ακμές διαφέρει.

\textbf{Σχόλιο: Αντίστροφη Διαγραφή.} Λιγότερο γνωστή αλλά κομψή παραλλαγή:
ξεκίνα με \textbf{όλες} τις ακμές, σάρωσε σε \textbf{φθίνουσα} σειρά κόστους, και
\textbf{διέγραψε} κάθε ακμή της οποίας η αφαίρεση δεν αποσυνδέει το γράφημα.
Λειτουργεί λόγω της ιδιότητας κύκλου. $O(m \log n)$ επίσης.

\section{Ιστορία των αλγορίθμων ΕΕΔ}

\begin{center}
\begin{tabular}{@{}l l@{}}
\toprule
Αλγόριθμος & Πολυπλοκότητα \\
\midrule
Borůvka (1926)                                & $O(m \log n)$ \\
Jarník (1930) / Prim (1957) / Dijkstra (1959) & $O(m \log n)$ \\
Kruskal (1956)                                & $O(m \log n)$ \\
Yao (1975) / Cheriton-Tarjan (1976)           & $O(m \log \log n)$ \\
Fredman-Tarjan (1987)                         & $O(m + n \log n)$ (Fibonacci heap) \\
Chazelle (2000)                               & $O(m \, \alpha(m, n))$ \\
Karger-Klein-Tarjan (1995)                    & $O(m)$ \textbf{τυχαιοποιημένος} \\
;                                             & $O(m)$ ντετερμινιστικά --- \textbf{ανοιχτό} \\
\bottomrule
\end{tabular}
\end{center}

Η ντετερμινιστική γραμμική πολυπλοκότητα παραμένει \textbf{ανοιχτό πρόβλημα}.
Μερικά από τα σημαντικότερα ονόματα της επιστήμης των υπολογιστών έχουν δουλέψει
σε αυτό.

\begin{recapbox}
\textbf{Τι κρατάμε από αυτή τη διάλεξη}
\begin{enumerate}
  \item \textbf{Dijkstra} --- άπληστος αλγόριθμος για συντομότερες διαδρομές με
  \textbf{θετικά} βάρη. Επεκτείνει την «εξερευνημένη ζώνη» κάθε φορά κατά την πιο
  κοντινή κορυφή. $O(m \log n)$ με σωρό.
  \item \textbf{ΕΕΔ} --- βρες δέντρο με $n - 1$ ακμές που συνδέει όλες τις
  κορυφές και έχει ελάχιστο συνολικό κόστος.
  \item \textbf{Ιδιότητα αποκοπής} --- η ελάχιστη ακμή σε κάθε αποκοπή ανήκει
  στο ΕΕΔ.
  \item \textbf{Ιδιότητα κύκλου} --- η μέγιστη ακμή σε κάθε κύκλο δεν ανήκει στο
  ΕΕΔ.
  \item \textbf{Prim} --- μεγαλώνει ένα δέντρο από μία κορυφή, διαλέγει σε κάθε
  βήμα τη φθηνότερη ακμή που το επεκτείνει. $O(m \log n)$.
  \item \textbf{Kruskal} --- σαρώνει ακμές σε αύξουσα σειρά κόστους, προσθέτει
  αν δεν δημιουργεί κύκλο. Χρησιμοποιεί Union-Find. $O(m \log n)$.
  \item \textbf{Στη χείριστη ανάλυση} οι τρεις αλγόριθμοι είναι ίδιοι. Στην
  εξεταστική, η εφαρμογή σε συγκεκριμένο γράφημα είναι ένα πολύ συχνό θέμα.
\end{enumerate}
\end{recapbox}

\lecturefooter
\end{document}
